\clearpage\SourceLeaf{136}
\begin{center}
[II.]

\medskip
\textbf{[INTERPOLATION DER COTANGENTEN UND COSECANTEN KLEINER BÖGEN.]}
\end{center}

Zur Interpolation der Cotangenten und Cosecanten kleiner Bögen sind
folgende Formeln die brauchbarsten:

Es gehören
\begin{center}
\begin{tabular}{cc}
zu den Bögen & die Cotangenten (oder Cosecanten) \\
$a$ & $\theta$ \\
$x$ & $y$ \\
$a'$ & $\theta'$, \\
\end{tabular}
\end{center}
so ist
\[
\begin{aligned}
\mathrm{I.}\quad x &= a+\frac{(\theta-y)a(a'-a)}{(\theta-y)a+(y-\theta')a'}
=a+\frac{(\theta-y)(a'-a)}{\theta-\theta'+\dfrac{a'-a}{a}(y-\theta')}\\[5pt]
&=a'-\frac{(y-\theta')a'(a'-a)}{(\theta-y)a+(y-\theta')a'}
=a'-\frac{(y-\theta')(a'-a)}{\theta-\theta'+\dfrac{a'-a}{a'}(y-\theta)}.
\end{aligned}
\]
\[
\begin{aligned}
\mathrm{II.}\quad y&=\frac{a'\theta'-a\theta}{a'-a}+\frac{(\theta-\theta')aa'}{(a'-a)x}
=\theta-\frac{(\theta-\theta')(x-a)a'}{(a'-a)x}\\[5pt]
&=\theta'+\frac{(\theta-\theta')(a'-x)a}{(a'-a)x}\,.
\end{aligned}
\]

Beispiel zu II.\quad Man sucht $\mathrm{cotg}\,0{,}060335$.
\begin{center}\small
\begin{tabular}{r@{\;}r@{\quad}l@{\;}l}
$\theta=$ & $10{,}5259499$ &&\\
$\theta'=$ & $10{,}5084181$ &&\\\cline{2-2}
$\theta-\theta'=$ & $175318$ & $\log\; .\ .\ .\ .\ .\ .\ .$ & $8{,}2438265-10$\\
$x-a=$ & $0{,}000035$ & $\log\; .\ .\ .\ .\ .\ .\ .$ & $5{,}5440680-10$\\
\multicolumn{2}{r}{$[a'=0{,}0604]$} & $\log\; .\ .\ .\ .\ .\ .\ .$ & $8{,}7810369-10$\\
&& $\mathrm{Compl.}\;\log x(a'-a)\; .\ .$ & $5{,}2194307$\\\cline{3-4}
\multicolumn{4}{r}{$0{,}00614274\; .\ .\ .\ .\ .\ .\ .\ .\ .\ .\ .\ .\ .\ .\ .\ .\ .\ .\ .\ .\;7{,}7883621-10$}\\
&$10{,}5259499$&&\\\cline{2-2}
&$10{,}5198072$&&
\end{tabular}
\end{center}
